\section{Contrastive Forecasting}
\label{sec:method}

Contrastive Forecasting consists of a per-patch encoder and a causal
forecaster, trained jointly with a contrastive loss in latent space.
The subsections below define the patching, the encoder, the
forecaster, and the loss in turn.

\subsection{Architecture}
\label{sec:architecture}

We work with a multivariate time series of shape $T \times C \times F$,
where $T$ is the number of raw timesteps, $C$ is the number of
channels, and $F$ is the number of features per timestep. In the
experiments reported here $F = 1$, but the formalism does not change
when $F > 1$. The series is cut into non-overlapping patches of width
$W$ along the time axis, so that the patched input has shape
$T_p \times C \times W \times F$ with $T_p = T / W$. Each patch is
encoded independently into a latent $h[t, c] \in \mathbb{R}^H$, and we
project it onto the Euclidean unit sphere by $L^2$-normalising, so
that $\|h[t, c]\| = 1$.

The choice of $W$ is dataset-dependent. Values in $\{16, 32, 64, 128\}$
all appear in the recent foundation-model literature. We use $W = 32$
in the experiments reported here, following the patch size used by
TimesFM~\citep{das2024timesfm}. To minimise compute, we set the stride between
consecutive patches to the patch width, giving non-overlapping windows
that cover the full sequence. A stride of $1$ is also of interest,
since it gives one latent per timestep, provided the loss is adjusted
to align each forecast with the correct future patch.

Our architecture has two models: an encoder and a forecaster. The
encoder maps a patch to a latent; the forecaster maps the sequence of
past latents to a forecast of the next latent. We write $f[t, c]$ for the forecaster's output at
position $t$ and channel $c$. The training objective is to make
$f[t, c]$ as close as possible to $h[t+1, c]$ under some dissimilarity
$D$,
\[
  D\big(f[t, c],\, h[t+1, c]\big) \to 0.
\]
Placing the latents on the unit sphere is a deliberate choice: it lets
us take $D$ to be the angular distance, measured through its proxy
the inner product (Section~\ref{sec:loss}). The representation also
benefits from living on a closed surface.

\subsection{Encoder}
\label{sec:encoder}

The first encoder we tried is the residual block used by
TiDE~\citep{das2023tide} and TimesFM~\citep{das2024timesfm}: a small
MLP (linear, ReLU, dropout, linear) summed with a linear skip from
the input, followed by a layer normalisation. The input dimension is
$W$ and the output dimension is $H$.

The encoder we currently use is a bidirectional GRU. Each patch is
read as a sequence of $W$ scalar timesteps, which gives the encoder a
chance to model the local temporal structure that an MLP collapses
into a single linear projection. The forward and backward final hidden
states ($128$ units each in our setup) are concatenated to a $256$-dim
vector and projected to $H$, summed with a linear skip from the raw
patch, and passed through a layer normalisation.

The encoder design space remains largely unexplored. In theory the
encoder could be as complex as a transformer, if encoding a patch into
a useful latent turns out to be a hard problem in its own right.
Finding the proper encoder is a potential future-work axis.

\subsection{Forecaster}
\label{sec:forecaster}

The forecaster is a causal transformer decoder. It has $L$ layers and
hidden width $H$, with $L$ and $H$ controlling the depth and the
capacity. Each layer applies a self-attention block, then a small
depthwise causal convolution (kernel size $3$) on the residual stream,
then a feed-forward network; residual connections and layer
normalisations are placed in the standard positions.

The convolution follows the observation of Elhage et
al.~\citep{elhage2021circuits} that a short-range depthwise
convolution lets induction-head-like circuits form without the second
attention layer that the canonical induction circuit requires. This
significantly improves the performance of small transformers, the
regime we operate in.


\subsection{Loss function}
\label{sec:loss}

For two unit-norm latents we score their similarity by their inner
product, which is the cosine of the angle between them and is
convenient to compute.

We index each latent and each forecast by the batch element
$b \in B$, the patch position $t \in T_p$, and the channel $c \in C$.
For every anchor $(b, t, c)$, the positive pair is the forecast
$f[b, t, c]$ paired with the future latent $h[b, t+1, c]$.

The negative set is built by varying anchors along three independent
axes: cross-temporal ($t \neq t'$), cross-channel ($c \neq c'$), and
cross-batch ($b \neq b'$). Each axis can be combined with the others,
and each pair can mix encoded latents $h$ and forecasts $f$. The loss
uses three specific pair families.
\begin{itemize}
\item \textbf{Cross-temporal} ($\mathcal{N}_{\mathrm{ct}}$): two
  encoded latents from the same series and channel at consecutive
  times, $\big(h[b, t, c],\, h[b, t+1, c]\big)$. Per anchor, $1$
  candidate.
\item \textbf{Cross-channel} ($\mathcal{N}_{\mathrm{cc}}$): two
  encoded latents from the same series at the same time, with
  different channels, $\big(h[b, t, c],\, h[b, t, c']\big)$ for
  $c' \neq c$. Per anchor, $C - 1$ candidates.
\item \textbf{Cross-batch} ($\mathcal{N}_{\mathrm{cb}}$): the future
  latent of one series paired with the forecast of a different series
  at the same prediction position,
  $\big(h[b, t+1, c],\, f[b', t, c]\big)$ for $b' \neq b$. Per anchor,
  $B - 1$ candidates.
\end{itemize}

Many other pairings are conceivable. The cross-temporal axis can also
be applied to $\big(h[b, t, c],\, f[b, t, c]\big)$ or
$\big(f[b, t, c],\, f[b, t+1, c]\big)$. The cross-batch axis can also
be applied to $\big(h[b, t, c],\, h[b', t, c]\big)$,
$\big(f[b, t, c],\, f[b', t, c]\big)$, or
$\big(h[b, t, c],\, f[b', t, c]\big)$. Any of these can be combined
with a channel difference. In general, adding more terms to the
denominator gives a faster loss decrease per backward pass at the cost
of slower loss computation, and the trade-off is not always
favourable; we specify the subset used in our experiments in
Section~\ref{sec:experiments}.

Cross-channel candidates are computed by a $C \times C$ inner-product
grid per batch element (diagonal masked), cross-batch candidates by a
$B \times B$ grid over the whole batch (diagonal masked), and
cross-temporal candidates by a single elementwise dot product over the
encoded tensor. At every position $(t, c)$, the candidates are aggregated
into three sums of size $B$ ($\mathcal{N}_{\mathrm{ct}}$), $B(C-1)$
($\mathcal{N}_{\mathrm{cc}}$), and $B(B-1)$
($\mathcal{N}_{\mathrm{cb}}$), and the same denominator is used for
every anchor at that position. The cross-batch set is therefore
quadratic in $B$ and dominates the negatives.

For unit-norm $u, v \in \mathbb{R}^H$, write the temperature-scaled
exponential of the inner product as
\[
  \sigma_\tau(u, v) \;=\; \exp\!\big(\langle u, v \rangle / \tau\big).
\]
The positive score at anchor $(b, t, c)$ is
\[
  \mathcal{P}_+(b, t, c) \;=\;
  \sigma_\tau\!\big(f[b, t, c],\, h[b, t+1, c]\big),
\]
and the three negative sums at position $(t, c)$ are
\begin{align*}
  \mathcal{N}_{\mathrm{ct}}(t, c)
    &\;=\; \sum_{b \in B}
      \sigma_\tau\!\big(h[b, t, c],\; h[b, t+1, c]\big),\\
  \mathcal{N}_{\mathrm{cc}}(t, c)
    &\;=\; \sum_{b \in B} \, \sum_{c' \neq c}
      \sigma_\tau\!\big(h[b, t, c],\; h[b, t, c']\big),\\
  \mathcal{N}_{\mathrm{cb}}(t, c)
    &\;=\; \sum_{b \neq b'}
      \sigma_\tau\!\big(h[b, t+1, c],\; f[b', t, c]\big).
\end{align*}
The contribution of the anchor $(b, t, c)$ to the loss is
\begin{equation}
  \mathcal{L}_{b, t, c}
  \;=\; -\log \frac{\mathcal{P}_+(b, t, c)}
  {\mathcal{N}_{\mathrm{ct}}(t, c) \;+\;
   \mathcal{N}_{\mathrm{cc}}(t, c) \;+\;
   \mathcal{N}_{\mathrm{cb}}(t, c)}.
  \label{eq:loss-anchor}
\end{equation}
The denominator depends on $(t, c)$ but not on $b$: every anchor at
the same position shares the same denominator. The mini-batch loss is
the average over all anchors,
\begin{equation}
  \mathcal{L} \;=\; \frac{1}{N}
  \sum_{b \in B} \sum_{t} \sum_{c \in C}
  \mathcal{L}_{b, t, c},
  \qquad N = B \cdot (T_p - 1) \cdot C.
  \label{eq:loss-batch}
\end{equation}
In the contrastive learning literature $\tau$ is usually chosen between
$0.05$ and $0.5$; the value used in our experiments is reported in
Section~\ref{sec:experiments}.
